\documentclass[11pt]{article}
\usepackage[margin=1in]{geometry}
\usepackage{amsmath,amssymb,amsthm,booktabs,hyperref}
\newtheorem{lemma}{Lemma}
\newtheorem{proposition}[lemma]{Proposition}
\newtheorem{theorem}[lemma]{Theorem}
\newcommand{\Word}[1]{W_{#1}}
\title{Transparent decoder routing and an exact bootstrap lock}
\author{Research note for the ProveIt Busy Beaver investigation}
\date{18 October 2026}
\begin{document}
\maketitle

\paragraph{Scope.}
This note concerns the old-label Prefix246 table with SHA-256
\texttt{303911c554284c49ce09042d8127feba4da260baae766e32260034aa3bdb82ee}.
The pinned repository is ProveIt commit
\texttt{c6bf052c70bbbafba8c39370a6a2d123f57cd2c1}.
It establishes local transformation lemmas and an exact finite bootstrap
invariant. A satisfying grammar is not, by itself, a complete machine
correctness certificate.

\section{Grammar states and transparent projections}

Let a semantic word be a nonempty word over the 44 context-labelled atomic
programs in the inherited grammar instance. A rule $w=(u,v)$ has $w=uv$ and
is implemented by one binary Turing-machine state with transitions
\[
  0\mapsto 0R\,\operatorname{entry}(u),\qquad
  1\mapsto 1R\,\operatorname{entry}(v).
\]
Every such transition preserves the tape cell it reads. Define the partial
projection $P_G(w,\epsilon)=w$ and
$P_G(w,b\alpha)=P_G(w_b,\alpha)$ when $w$ has the rule
$(w_0,w_1)$. A projection is undefined if it enters an atom before exhausting
its path. Distinct semantic words use distinct grammar states; context-labelled
atoms may share a physical state only as prescribed by the input instance.

The words are sequences of \emph{closed} atomic programs. This qualification
matters: the underlying register language has structured break and continue
operations, so arbitrary tree association is not automatically sound. For a
tree $t$, the static escape bounds are
\[
 b(\operatorname{node}(u,v))=\max(b(u)-1,b(v),0),\qquad
 c(\operatorname{node}(u,v))=\max(c(u),c(v)-1,0).
\]
A decrement leaf has break bound~1; a break-$k$ or continue-$k$ leaf has the
corresponding bound~$k$; other bounds are zero. Direct recomputation gives
$(b,c)=(0,0)$ for all 44 atoms and both children of each of the 127 original
sequence states. Accordingly, their ordinary sequential composition has the
usual associativity law. This justifies the abstract word algebra; it does not
replace the tape-level decoder checks.

\begin{lemma}[Forced-prefix routing]
Suppose a caller at head position $h$ reads $b$, writes $d$, moves left,
and enters a transparent decoder. Let $x$ be the arbitrary tape bit at $h-1$.
Suppose its old route follows a word $\alpha_x$ of length one or two, where
$\alpha_x$ is the corresponding prefix of $xd$, and ends at a specified
semantic word or atom $v_x$.

Replace the caller's target by an active grammar word $w$ such that
$P_G(w,\alpha_x)=v_x$ for both $x=0,1$. Then the new and old routes take
exactly $1+|\alpha_x|$ steps, have the same tape after the block, and end at
head position $h-1+|\alpha_x|$. Their endpoint states represent the same
specified word or context-labelled atom. If that endpoint retains its old
state label, the two endpoint configurations are literally identical.
\end{lemma}

\begin{proof}
After the caller's step, the head is at $h-1$, the bit there is $x$, and the
bit at $h$ is $d$. Each subsequent decoder step preserves its current bit
and moves one cell right. Thus the only bits read are the specified prefix
of $xd$; no assumption about the remainder of the tape is used. The
projection equations give the endpoint after exactly the stipulated number
of steps. The statement about tape, head position and retained endpoint
labels follows directly.
\end{proof}

\section{Removing the fixed routing role of state 177}

The literal table has precisely the following incoming edges to state~177:
\[
 153:\;1\mapsto 1L177,\qquad
 159:\;0\mapsto 0L177.
\]
Its two outgoing edges are $0\mapsto0R175$ and $1\mapsto1R207$.
The relevant children are
\[
 175=(113,105),\qquad 207=(213,231).
\]
Here $\Word q$ denotes the semantic word of old sequence state~$q$.

\begin{proposition}[Eliminating a fixed router]
Choose active grammar words $U,V$ satisfying
\[
\begin{array}{ll}
 P_G(U,01)=\Word{105}, & P_G(U,11)=\Word{231},\\
 P_G(V,00)=\Word{113}, & P_G(V,10)=\Word{213}.
\end{array}
\]
Retarget $153/1$ to $U$ and $159/0$ to $V$, retaining their writes and
leftward moves. The fixed routing implementation at state~177 is then
unnecessary. Each replaced block has three steps and satisfies the preceding
lemma on every tape. State~177 has no other incoming edge.
\end{proposition}

The original table has 127 sequence states and 119 remaining working states.
After this change, a grammar with $K$ active rules therefore costs
$118+K$ working states, provided every unchanged external entry has an active
target and all erased incoming edges are discharged. In particular, $K=127$
would give 245 states. A new grammar state may reuse numeric label~177;
the saving concerns its fixed implementation role, not a prohibition on
reusing its label. The original external entry requirements at~175 and~207
come only from~177; their other incoming edges are from replaceable sequence
states.

\section{A bootstrap lock that leaves 100 sequence states unfixed}

Two independent exact interpreters---a sparse set tape and a dense array
tape with array-indexed transition rows---were run for $410735$ steps from
the blank tape and start state~95. They agree on the endpoint, every
surviving one-cell, and every used transition count. The complete derivation
is reproducible by \texttt{generate\_bootstrap\_lock.py}; its output is
\texttt{bootstrap\_lock.json}.

The trace uses 55 distinct sequence rows from 28 sequence states. Its only
uses of caller rows eligible for the transparent rerouting transformation are
\[
\begin{array}{c|c|c|r}
\text{caller/read}&x&\alpha_x&\text{occurrences}\\\hline
78/0&1&10&6\\
92/1&0&0&12\\
92/1&1&1&3\\
151/1&0&0&21
\end{array}
\]
The endpoints are respectively old states~123, 36, 37 and~36. State~177
and its two incoming caller rows are not used in this bootstrap.

Treat each listed caller occurrence together with its one or two following
transparent transitions as a block. Excluding these block interiors leaves
54 distinct sequence rows from 27 states: only old row $89/1$ disappears
from the set of distinct required rows. These 54 rows comprise both rows of
each of 27 original grammar states. Thus 100 of the original 127 sequence
states are not directly fixed by this lock.

\begin{theorem}[Exact block-aware bootstrap lock]
Consider an assembled grammar machine satisfying all of the following:
every unchanged framework row is retained; every new grammar state has its
transparent two-row implementation; all rerouting blocks used by the
baseline bootstrap satisfy the forced-prefix lemma; and, for each of the
54 listed sequence rows $(q,b)$, the word $\Word q$ is active with
\[
 P_G(\Word q,b)=\text{the original semantic child word of row }(q,b).
\]
Allocate every retained original semantic word to its original state label,
and allocate new words injectively without colliding with unchanged states.
Then the new machine has exactly the baseline configuration after $410735$
steps from the blank tape.
\end{theorem}

\begin{proof}
Partition the independently replayed baseline trace into the stated
rerouting blocks and singleton steps outside those blocks. At a singleton
framework step the row is unchanged. At a singleton sequence step the lock
forces the same original child word. Rule closure forces any non-atomic
child to remain active, and label-preserving allocation therefore forces
the same literal successor state; an atomic child already uses its fixed
physical state. Thus this singleton step preserves equality of complete
configurations.

At a rerouting block the lemma preserves tape, head position and duration.
The actual endpoints~123, 36 and~37 retain their old state labels: state~123
is among the locked sequence states, while~36 and~37 are unchanged atomic
states. Thus equality holds again at the block's end. Induction over the
finite partition proves equality at its final cut. The trace ends at the
dispatcher, outside every block interior, so the final cut is exactly step
$410735$ in both machines.
\end{proof}

The resulting configuration has dispatcher state~96, head~$-5$, marker~$-6$,
header~176, an empty instruction pointer and header distance~181. The nine
physical registers are
\[
 (0,0,0,3,0,0,3,0,0),
\]
and the logical registers $(N,T,W,A,D,H,X,Y,L)$ are
$(3,0,0,0,3,0,0,0,0)$. There are exactly 167 spare blocks, all of value~1,
and 351 one-cells in total; every other cell is blank. This supplies the
same concrete root seed as the reviewed Prefix246 construction.

\section{An optional shortened scanner}

The framework contains the transparent-on-zero block
\[
 58:\;0\mapsto0L74,\qquad
 74:\;0\mapsto0L84,\qquad
 84:\;0\mapsto0R84.
\]
If the three visited cells are zero, these three steps can be replaced by
the single transition $58:\;0\mapsto0L84$: tape, final head position and
final state are identical. State~74 has only the incoming edge $58/0$, so
its fixed scanner role can then be removed.

The fresh baseline replay has exactly 1410 occurrences of this three-step
block; in every occurrence, all three required cells are zero. It never
uses row $74/1$. The generator also independently replays the shortened
table with a dense tape for exactly
\[
 410735-2\cdot1410=407915
\]
steps and obtains precisely the same complete root frame.

\begin{proposition}[Bootstrap lock with scanner shortening]
Under the block-aware bootstrap theorem's conditions, additionally remove
the fixed scanner state~74 and replace row $58/0$ by $0L84$. Then the
same complete root frame is reached at step~407915. The working-state
count is $118+K$ without the router-177 elimination, or $117+K$ with it.
\end{proposition}

\begin{proof}
Include the 1410 guarded three-step scanner occurrences in the finite trace
partition used above. They do not overlap the decoder routing blocks.
Each has the same complete endpoint after one new step instead of three
old steps. Induction over the partition therefore preserves the endpoint
and subtracts exactly $2\cdot1410$ from its time. Both state-count equations
follow by counting the removed fixed implementation roles.
\end{proof}

This is an exact finite \emph{bootstrap} statement. The shorter scanner is
known to fail some instructions of the original source. A proposed repair
must reduce the source zero-run bound and pass every guarded instruction
simulation against the \emph{actual shortened table}. A six-zero grammar
bound is a search constraint, not a substitute for that runtime proof.

\paragraph{Remaining obligations for a new small machine.}
The bootstrap lock protects the blank start, and the routing equations
protect the affected local decoder entries. The emitted whole machine must
still pass the complete register-program correspondence and all instruction
simulation checks, including the zero-run scanner condition and the macro
contracts used by symbolic replay. Finite-pool unsatisfiability is only an
obstruction within the specified search family. It is not a lower bound
on the number of states in machines deciding the target consistency problem.

\end{document}
